\documentclass[11pt]{article}

\usepackage[margin=1in]{geometry}
\usepackage{amsmath,amssymb,amsthm}
\usepackage[hidelinks]{hyperref}

\newtheorem{theorem}{Theorem}
\newtheorem{lemma}{Lemma}
\newtheorem{corollary}{Corollary}

\newcommand{\N}{\mathbb{N}}
\newcommand{\K}{\mathbb{K}}
\newcommand{\X}{\mathcal{X}}
\newcommand{\Exc}{\operatorname{p\text{-}Exc}}
\newcommand{\spancl}{\overline{\operatorname{span}}}
\newcommand{\rank}{\operatorname{rank}}

\title{Finite-Dimensional p-Excess Invariance for Approximately Dual p-ASFs}
\author{agent\_lane\_02}
\date{June 18, 2026}

\begin{document}
\maketitle

\section*{Source and Question}

The source paper is K. Mahesh Krishna and P. Sam Johnson,
\emph{Approximately Dual p-Approximate Schauder Frames},
arXiv:2110.10121.  Problem 3.3 asks whether, for two approximately
dual p-approximate Schauder frames
\[
  (\{f_n\}_{n},\{\tau_n\}_{n}), \qquad
  (\{g_n\}_{n},\{\omega_n\}_{n})
\]
for a Banach space \(\X\), one always has
\[
  \Exc(\{f_n\}_{n},\{\tau_n\}_{n})
  =
  \Exc(\{g_n\}_{n},\{\omega_n\}_{n}).
\]
Here
\[
\Exc(\{f_n\},\{\tau_n\})
=
\sup\Bigl\{|M|:M\subseteq\N,\ 
\spancl\{\tau_n:n\notin M\}=\X,\ 
\spancl\{f_n:n\notin M\}=\X^*
\Bigr\}.
\]

\section*{Proof Intuition}

In finite dimension, a p-ASF is already controlled by arbitrarily large
finite chunks.  Since the analysis map has finite-dimensional domain, the
coordinate truncations converge uniformly on its range in \(\ell^p\).  Thus
large partial frame operators are invertible.  For a finite invertible
partial frame operator, Cauchy-Binet forces at least one \(d\)-element set of
indices whose vectors span \(\X\) and whose functionals span \(\X^*\).  All
other indices in that finite chunk may be deleted while those \(d\) indices
are left behind.  Letting the chunk grow gives infinite p-excess.

\section*{Finite Matrix Lemma}

\begin{lemma}\label{lem:finite}
Let \(\dim \X=d<\infty\), and let
\((f_i,\tau_i)_{i=1}^N\) be a finite p-ASF analogue on \(\X\), meaning that
the operator
\[
  Sx=\sum_{i=1}^N f_i(x)\tau_i
\]
is invertible.  Then
\[
  \Exc((f_i)_{i=1}^N,(\tau_i)_{i=1}^N)=N-d.
\]
\end{lemma}

\begin{proof}
No deletion set can have cardinality larger than \(N-d\), since fewer than
\(d\) remaining vectors cannot span \(\X\), and fewer than \(d\) remaining
functionals cannot span \(\X^*\).

Choose a basis of \(\X\) and the dual coordinate basis of \(\X^*\).  Put the
vectors \(\tau_i\) as the columns of a \(d\times N\) matrix \(T\), and put
the functionals \(f_i\) as the rows of an \(N\times d\) matrix \(F\).  The
matrix of \(S\) is \(TF\), hence \(\det(TF)\ne0\).  By Cauchy-Binet,
\[
  \det(TF)
  =
  \sum_{\substack{I\subseteq\{1,\ldots,N\}\\ |I|=d}}
     \det(T_I)\det(F_I),
\]
where \(T_I\) is the \(d\times d\) matrix of columns indexed by \(I\), and
\(F_I\) is the \(d\times d\) matrix of rows indexed by \(I\).  Therefore
some \(I\) has both determinants nonzero.  The vectors
\(\{\tau_i:i\in I\}\) span \(\X\), and the functionals
\(\{f_i:i\in I\}\) span \(\X^*\).  Deleting the complementary set
\(\{1,\ldots,N\}\setminus I\), of cardinality \(N-d\), is therefore allowed.
This proves the reverse inequality.
\end{proof}

\section*{Finite-Dimensional Countable Case}

\begin{theorem}\label{thm:finite-dim}
Let \(1\le p<\infty\), let \(\dim \X=d<\infty\), and let
\((\{f_n\}_{n=1}^\infty,\{\tau_n\}_{n=1}^\infty)\) be a p-ASF for \(\X\).
Then
\[
  \Exc(\{f_n\},\{\tau_n\})=\infty .
\]
\end{theorem}

\begin{proof}
Let \(\theta_f:\X\to\ell^p\) and \(\theta_\tau:\ell^p\to\X\) be the
analysis and synthesis operators.  The frame operator
\[
  S=\theta_\tau\theta_f
\]
is invertible.  Let \(P_N:\ell^p\to\ell^p\) be the coordinate projection
onto the first \(N\) coordinates and set
\[
  S_N=\theta_\tau P_N\theta_f
      =\sum_{n=1}^N f_n(\,\cdot\,)\tau_n .
\]
Because \(\X\) is finite dimensional, the image of its unit ball under
\(\theta_f\) is compact in \(\ell^p\).  Since \(P_Ny\to y\) for each
\(y\in\ell^p\), the convergence is uniform on this compact image.  Hence
\[
  \|S_N-S\|
  \le \|\theta_\tau\|\,\|(P_N-I)\theta_f\|
  \longrightarrow 0 .
\]
For all sufficiently large \(N\), \(S_N\) is invertible.  Applying
Lemma~\ref{lem:finite} to the finite family
\((f_n,\tau_n)_{n=1}^N\), we obtain a set \(I_N\subseteq\{1,\ldots,N\}\)
with \(|I_N|=d\) such that \(\{\tau_n:n\in I_N\}\) spans \(\X\) and
\(\{f_n:n\in I_N\}\) spans \(\X^*\).  Therefore the deletion set
\[
  M_N=\{1,\ldots,N\}\setminus I_N
\]
is admissible for the original countable p-ASF, because the remaining
countable family still contains the spanning subfamilies indexed by \(I_N\).
Thus
\[
  \Exc(\{f_n\},\{\tau_n\})\ge |M_N|=N-d
\]
for arbitrarily large \(N\), and the p-excess is infinite.
\end{proof}

\begin{corollary}
Problem 3.3 has an affirmative answer for finite-dimensional Banach spaces.
Indeed, if \(\dim \X<\infty\), any two p-ASFs for \(\X\), approximately dual
or not, have the same p-excess, namely \(\infty\).  For the finite-index
analogue with \(N\) atoms, any two finite p-ASFs for \(\X\) have p-excess
\(N-\dim\X\).
\end{corollary}

\section*{Limitations}

The proof does not settle the infinite-dimensional case of Problem 3.3.  In
infinite dimension, finite coordinate truncations need not approximate the
frame operator in operator norm, and the Cauchy-Binet common-subset argument
has no direct infinite-dimensional replacement.  The additional requirement
that the remaining functionals span \(\X^*\), not only that the remaining
vectors span \(\X\), is the point that still requires a genuinely
infinite-dimensional argument or a counterexample.

\section*{Novelty Check}

Local run indexes and attempts contained no existing result packet for
arXiv:2110.10121 or Problem 3.3.  Web searches used the phrases
\begin{itemize}
\item \texttt{"p-Exc" "approximately dual"},
\item \texttt{"Approximately Dual p-Approximate Schauder Frames" "excess"},
\item \texttt{"p-excess" "p-approximate Schauder frames"}.
\end{itemize}
These searches found the source paper and the Hilbert-frame excess paper of
Bakic and Beric, but no later paper resolving the p-ASF problem.  The
finite-dimensional result above is therefore recorded as a run-generated
partial result, subject to human review.

\begin{thebibliography}{9}

\bibitem{KJ}
K. Mahesh Krishna and P. Sam Johnson,
\emph{Approximately Dual p-Approximate Schauder Frames},
arXiv:2110.10121, 2021.

\bibitem{BB}
Damir Bakic and Tomislav Beric,
\emph{On excesses of frames},
Glasnik Matematicki Serija III 50(70) (2015), no. 2, 415--427;
arXiv:1403.0632.

\end{thebibliography}

\end{document}
